% Escolhas editoriais deste projeto, e não medidas impostas pela revista:
% A4; margens de 2,5 cm; corpo 12; entrelinhas 1,5; notas 10, simples.
\usepackage[T1]{fontenc}
\usepackage[utf8]{inputenc}
\usepackage[english,brazil]{babel}
\usepackage{lmodern}
\usepackage[margin=2.5cm]{geometry}
\usepackage{setspace}
\usepackage{microtype}
\usepackage[hang,flushmargin,bottom]{footmisc}
\usepackage{needspace}
\usepackage{titlesec}
\usepackage{etoolbox}
\usepackage{xurl}
\usepackage[hidelinks,unicode]{hyperref}
\usepackage{csquotes}
\hypersetup{pdfauthor={},pdfsubject={},pdfkeywords={},pdftitle={}}
\onehalfspacing
\setlength{\parindent}{1.25cm}
\setlength{\parskip}{0pt}
\setlength{\footnotemargin}{1.5em}
\setlength{\footnotesep}{8pt}
\setlength{\skip\footins}{14pt plus 3pt minus 2pt}
\renewcommand{\footnotesize}{\fontsize{10}{12}\selectfont}
% Notas extensas podem continuar na página seguinte, em vez de transbordar.
% Um valor maior procura manter cada nota unida; não use 10000 para notas longas.
\interfootnotelinepenalty=100
\setcounter{secnumdepth}{0}
\titleformat{\section}{\normalfont\normalsize\bfseries}{}{0pt}{}
\titleformat{\subsection}{\normalfont\normalsize\bfseries}{}{0pt}{}
\titlespacing*{\section}{0pt}{1.1\baselineskip}{0.5\baselineskip}
\titlespacing{\subsection}{0pt}{0.9\baselineskip}{0.35\baselineskip}
\makeatletter
\renewcommand{\maketitle}{%
  \begin{center}
  {\fontsize{16}{20}\selectfont\bfseries\@title\par}
  \end{center}
  \vspace{0.5\baselineskip}}
\makeatother
\AtBeginEnvironment{quote}{\singlespacing\small}
\AtBeginEnvironment{quotation}{\singlespacing\small}
\AtBeginEnvironment{description}{\singlespacing}
\pretocmd{\subsection}{\Needspace{5\baselineskip}}{}{}
\emergencystretch=1.5em
\widowpenalty=1000
\clubpenalty=1000
\urlstyle{same}
